\section{结构控制、视频与下一代生成架构}

图像生成不止文本条件。Pose、edge、scribble 等结构信号具有空间对应，subject-driven generation 与 editing 则可能没有直接对齐；视频进一步增加时间维度。最后几页把 DiT 从单图 backbone 扩展为更广的 conditional generation 平台。

\subsection{结构控制不是统一问题}

若控制信号与输出像素位置对应，可用 auxiliary encoder 产生多尺度 features 并注入 base model；若条件是主体身份或编辑指令，位置关系不固定，需要 token interaction、adapter 或检索式机制。本节从“条件是否具有空间 correspondence”建立分类，因为这决定条件能否逐位置注入，也决定评估应侧重轮廓遵循、身份保持还是编辑语义。

\lecturefigure{slide-59.jpg}{结构控制示例：edge、depth、pose 与目标图像具有空间对应}{59}

这些信号可视为与 noisy latent 对齐的额外观测。模型应保持主体内容，同时严格遵循轮廓或姿态。评估不能只看美观，还需计算 control accuracy 与 content preservation。

\lecturefigure{slide-60.jpg}{ControlNet/PixArt-$\delta$：辅助网络编码视觉条件并注入 base DiT}{60}

Auxiliary network 复制或适配部分 backbone，把结构 features 通过 zero-initialized residual 注入。这样可冻结大 base model，减少微调成本。对于无空间对应的 subject/edit conditions，直接逐位置相加并不合理。

\teachervoice{01:08:47--01:10:40，老师特别区分结构控制与 subject-driven/editing：前者总有空间 correspondence，后者未必有，因此不能把 ControlNet 注入方式机械推广。}

\begin{warningbox}{控制强度与生成自由度冲突}
过强结构条件会锁死纹理和构图，过弱又无法遵循控制。应报告不同 control scale 下的 fidelity--diversity 曲线，并检查条件噪声与域外姿态。
\end{warningbox}

\subsection{视频：从二维 token 到三维时空交互}

视频不仅多了帧数，还要求跨时间保持身份、运动和物理一致。把 $T$ 帧、每帧 $N$ 个 token 全部 joint attention，成本为 $O((TN)^2)$；factorized attention 更省，却可能损失长程时空耦合。本节把单图 DiT 的 token accounting 扩展到时间轴，重点比较 full 3D interaction、空间--时间分解和位置编码分别解决什么问题。

\lecturefigure{slide-61.jpg}{视频 DiT：RoPE、3D attention、factorized attention 与效率挑战}{61}

RoPE（Rotary Position Embedding）通过旋转 query/key 编码相对位置，可扩展到时间与空间轴。Full 3D attention 质量通常更强但成本极高；空间/时间分解 attention 更可控。视频指标还需覆盖 temporal consistency，而非逐帧清晰度。

\teachervoice{00:59:10--01:01:38 与 01:10:40--01:11:44，问答强调视频的困难不是“逐帧多生成几张图”，而是三维注意力成本与时序一致性；老师建议关注 LTX-Video 等效率路线。}

\begin{knowledgebox}{第一次出现：RoPE}
RoPE 把位置编码为 query/key 的旋转，使点积自然包含相对位置信息。视频可为时间、高度、宽度分配旋转频率，但长度外推和三轴比例仍需设计。
\end{knowledgebox}

\subsection{In-context generation 与开放方向}

下一代生成模型希望像语言模型一样在上下文中读取示例：给参考图、编辑前后对或多模态序列，模型直接推断任务。现有架构的主要障碍是序列太长、条件类型异构以及训练数据难构造。本节从固定接口的 conditional generation 走向示例驱动任务，阅读后面几页时要区分“统一输入格式”的演示与“真正跨任务泛化”的证据。

\lecturefigure{slide-62.jpg}{下一代问题：如何在扩散模型中实现 in-context learning}{62}

In-context generation 需要把示例与目标放进同一可交互上下文，而不为每个任务微调。若所有图像都展开为 tokens，计算迅速爆炸；若过度压缩，又无法复制细节或关系。架构必须结合高效 attention、分层表示与任务 token。

\lecturefigure{slide-63.jpg}{Playground v3、FuseDiT、OmniGen 等探索统一的上下文生成}{63}

这些模型尝试把 text、reference images 与 target latents 组织成统一序列或分阶段交互。它们展示方向，而非已经解决通用 in-context learning。比较时应检查支持的任务、上下文长度、训练数据与是否需要专用 adapter。

\lecturefigure{slide-64.jpg}{Diffusers 提供多种架构实现，便于复现与 hacking}{64}

实现页强调研究入口：统一库可比较 DiT、PixArt、SANA、MMDiT 等模块，减少隐藏实现差异。但“代码可运行”不等于实验可复现；还需模型权重、数据处理、训练配置、随机种子和评估脚本。

\lecturefigure{slide-65.jpg}{结论与未覆盖方向：MoE、训练、post-training、alignment 与评估}{65}

最后一页主动列出空白：专家混合（MoE）已进入图像生成；训练本身是另一套系统；后训练、偏好对齐与评估仍是独立难题。这个边界阻止我们把 backbone 设计误写成完整生成智能。

\teachervoice{01:11:44--01:14:23，老师把 in-context generation 视为下一代架构问题，并明确承认没有覆盖 MoE、训练、post-training、preference alignment 与 evaluation。}

\begin{importantbox}{架构研究的终点不是一张更漂亮的 sample}
完整生成系统还要回答：数据是否授权与去重？训练是否稳定？模型如何对齐偏好？自动指标是否可信？视频是否时序一致？部署是否满足延迟与显存预算？Backbone 只是其中一层。
\end{importantbox}

\subsection{本章小结}

结构控制依赖条件是否空间对齐；视频引入三维 token 与时序一致性；in-context generation 则要求模型统一处理示例与目标。Diffusers 降低实现门槛，但训练、post-training、MoE 和评估仍是开放系统问题。

